Entity matching prototype — observations

Setup
- Test set: 20 messy company names covering 7 case types (suffixes,
  typos, reordered words, abbreviations, extra words, French names,
  and decoys that look like a different company).
- Matched against 10 recognized companies using rapidfuzz's
  process.extractOne, after normalizing both sides (lowercase, remove
  accents and punctuation, "&" -> "and", strip legal suffixes).

Results (correct out of 20)
- ratio: 15        partial_ratio: 18
- token_sort_ratio: 15   token_set_ratio: 18
- Best precision/recall: partial_ratio at threshold 90
  (precision 1.00, recall 0.80).

What I noticed
1. Shared filler words cause wrong matches. With plain ratio,
   "IBM Canada Ltd." matched Bell Canada because "canada" is longer
   than "ibm" and dominates the score.
2. Each scorer fixes some cases and breaks others. token_sort fixed
   reordered names ("Hydro Toronto") but broke "Hydro One Networks."
3. token_set's worst error can't be fixed with a threshold.
   "Bell Textron Canada" matched Bell Canada with a score of 100,
   because every word of "bell canada" is in it. Its precision stays
   at 0.94 at every threshold.
4. partial_ratio scored best here, but short names like "cgi" or
   "bell" can fit inside unrelated longer names, so I'd expect more
   false matches on a real list of thousands of suppliers.
5. My normalization turned "I.B.M." into "i b m". It still matched
   IBM, but with low scores (50-75), so a strict threshold sends it
   to review.
6. "International Business Machines" failed with every scorer.
   String similarity can't connect a full name to its abbreviation.

What I'd try next
- An alias table for known abbreviations (IBM, CGI).
- Merging single letters after normalization ("i b m" -> "ibm").
- Flagging matches where the input has extra distinguishing words
  (like "textron") for human review instead of auto-accepting.
- For government data: a false merge is worse than a missed match,
  so I'd favor precision and send uncertain cases to review.